% !TEX root = 抽象代数前两章全部习题详解.tex
\clearpage\section{第1章：代数基本概念}

\subsection{习题1：\texorpdfstring{平方公式判别交换群}{平方公式判别交换群}}
\label{ex:1:1}

\paragraph{题目}
若群$G$中任意$a,b$满足$(ab)^2=a^2b^2$，证明$G$交换。

\paragraph{考点与思路}
将幂展开，只消去两端的因子。也可证明该条件与交换性等价。

\paragraph{所用结论与原文依据}
\noteref{cancel}。

% 原解来源：chapters/Part 2/chapter1-exercises.tex:6；第1章原\exref{1}{1}。

\begin{proof}
任取$a,b\in G$，题设给出$abab=aabb$。先在两端左乘$a^{-1}$，再右乘$b^{-1}$，得到
\begin{align*}
abab&=aabb,\\
bab&=abb,\\
ba&=ab.
\end{align*}
因为$a,b$任意，群中任意两元素交换，故$G$为交换群。

反过来，若$G$交换，则$abab=aabb=a^2b^2$，所以题设条件也成立。
\end{proof}

\paragraph{复习提示}
消去律仅允许消去位于同一侧的相同因子；不能从一开始就把中间的$ba$换成$ab$。


\subsection{习题2：\texorpdfstring{每个元素的平方为单位元}{每个元素的平方为单位元}}
\label{ex:1:2}

\paragraph{题目}
若群$G$中的每个元素$a$都满足$a^2=e$，证明$G$交换。

\paragraph{考点与思路}
将题设同时用于$a$、$b$和$ab$，再使用乘积取逆时的反序公式。

\paragraph{所用结论与原文依据}
\noteref{glaws}。

% 原解来源：chapters/Part 2/chapter1-exercises.tex:14；第1章原\exref{1}{2}。

\begin{proof}
由$a^2=e$知$a$本身就是$a$的逆元，所以每个$a\in G$都有$a^{-1}=a$。
任取$a,b\in G$，因为$ab\in G$，它也满足$(ab)^{-1}=ab$。另一方面，任何群中都有$(ab)^{-1}=b^{-1}a^{-1}$，于是
\[
ab=(ab)^{-1}=b^{-1}a^{-1}=ba.
\]
因此$G$交换。
\end{proof}

\paragraph{复习提示}
题设包含单位元，并不意味着每个元素的阶都恰为2；单位元的阶为1。结论也不要求$G$有限。


\subsection{习题3：\texorpdfstring{有限消去半群是群}{有限消去半群是群}}
\label{ex:1:3}

\paragraph{题目}
非空有限集合上的结合运算若同时满足左右消去律，则构成群。

\paragraph{考点与思路}
先用有限集上单射等价于满射，构造单侧单位元，再把它们合并为双侧单位元。

\paragraph{所用结论与原文依据}
\noteref{group}；\noteref{glaws}。

% 原解来源：chapters/Part 2/chapter1-exercises.tex:21；第1章原\exref{1}{3}。

\begin{proof}
记给定集合为$G$。固定$a\in G$，由左右消去律，映射$x\mapsto ax$及$x\mapsto xa$均为$G$到自身的单射。$G$有限，故这两个映射也满射，因而存在$e,f\in G$满足$ae=a$、$fa=a$。

任取$x\in G$，结合律给出
\begin{align*}
a(ex)&=(ae)x=ax,\\
(xf)a&=x(fa)=xa.
\end{align*}
分别使用左、右消去律，得$ex=x$及$xf=x$。所以$e$是整个集合的左单位元，$f$是右单位元。于是$ef=f$且$ef=e$，从而$e=f$为双侧单位元。

对任意$g\in G$，左右乘$g$的映射满射，故可取$b,c\in G$使$gb=e$、$cg=e$。由
\[
c=ce=c(gb)=(cg)b=eb=b
\]
可知$gb=bg=e$。每个元素都有双侧逆元，故$G$为群。

有限性不能删去：正整数集在加法下结合且满足两侧消去律，但没有加法单位元，故不是群。
\end{proof}

\paragraph{复习提示}
尚未构造出群时，只能使用题设的消去律和结合律，不能预先使用逆元。


\subsection{习题4：\texorpdfstring{可解方程的结合系统}{可解方程的结合系统}}
\label{ex:1:4}

\paragraph{题目}
非空集合$G$上有一个结合的二元运算。若对任意$a,b\in G$，方程$ax=b$与$ya=b$都有解，证明$G$构成群。

\paragraph{考点与思路}
尚未证明是群时，不能使用逆元和消去律。先将某个元素的单侧单位元推广为全体元素的单侧单位元。

\paragraph{所用结论与原文依据}
\noteref{group}。

% 原解来源：chapters/Part 2/chapter1-exercises.tex:33；第1章原\exref{1}{4}。

\begin{proof}
\textbf{第一步：构造双边单位元。}
因为$G$非空，可固定$a\in G$。由题设，存在$e,f\in G$使
\[
ae=a,\qquad fa=a.
\]
任取$x\in G$。方程$ya=x$有解，因此$x=ya$，结合律给出
\[
xe=(ya)e=y(ae)=ya=x.
\]
所以$e$是整个集合的右单位元。又由方程$az=x$有解，可写$x=az$，从而
\[
fx=f(az)=(fa)z=az=x.
\]
所以$f$是整个集合的左单位元。将两者用于$fe$，分别得到$fe=f$与$fe=e$，故$f=e$。因此$e$为双边单位元。

\textbf{第二步：构造双边逆元。}
任取$g\in G$，由方程可解性，存在$b,c\in G$使$gb=e$、$cg=e$。于是
\[
c=ce=c(gb)=(cg)b=eb=b.
\]
所以$gb=bg=e$，即$g$有双边逆元。

二元运算已经具有封闭性，结合律由题设保证；现在单位元和每个元素的逆元也已得到，故$G$是群。
\end{proof}

\paragraph{复习提示}
“有解”不必加强为“唯一解”。在构造出群以前，任何一步都不应预先使用消去律。


\subsection{习题5：\texorpdfstring{非交换平方公式的反例}{非交换平方公式的反例}}
\label{ex:1:5}

\paragraph{题目}
举例说明群中一般不成立$(xy)^2=x^2y^2$。

\paragraph{考点与思路}
沿原解选取两个相邻对换，分别计算两个表达式。

\paragraph{所用结论与原文依据}
\noteref{cycles}。

% 原解来源：chapters/Part 2/chapter1-exercises.tex:44；第1章原\exref{1}{5}。

\begin{proof}
在$S_3$中取$x=(12)$、$y=(23)$。它们都是对换，所以$x^2=y^2=e$，从而$x^2y^2=e$。

按右边先作用的约定，$xy$将$1,2,3$依次映为$2,3,1$，故$xy=(123)$。因此
\[
(xy)^2=(132)\ne e=x^2y^2.
\]
这说明一般群中平方公式$(xy)^2=x^2y^2$不成立。
\end{proof}

\paragraph{复习提示}
计算置换乘积时始终从右边开始作用。


\subsection{习题6：\texorpdfstring{任意偶数阶的非交换例子}{任意偶数阶的非交换例子}}
\label{ex:1:6}

\paragraph{题目}
对整数$n>2$，构造一个$2n$阶的非交换群。

\paragraph{考点与思路}
沿用原解答，取正$n$边形的几何对称群，再按正文中相邻顶点像的计数和旋转、反射关系证明其阶数与非交换性。

\paragraph{所用结论与原文依据}
\noteref{dihedral}。

% 原解来源：chapters/Part 2/chapter1-exercises.tex:51；第1章原\exref{1}{6}。

\begin{proof}
取平面上的正$n$边形，记$D_n$为保持此正多边形整体不变的所有平面等距变换的集合。恒等变换属于$D_n$；两个对称变换的复合仍保持图形；逆变换也保持图形。映射复合具有结合律，所以$D_n$在复合下成群。

\textbf{先证明$|D_n|=2n$。}任意对称变换置换顶点，因而保持顶点的重心，即正多边形的中心$O$。固定两个相邻顶点$P,Q$，由于$n>2$，$O,P,Q$不共线。平面等距变换由三个不共线点的像唯一决定，因此一个对称变换由$P,Q$的像唯一决定。
$P$的像有$n$种选择；$Q$的像必须与$P$的像相邻，而每个顶点恰有两个相邻顶点，所以$Q$的像至多有两种选择。故$|D_n|\leq2n$。

令$T$为绕$O$旋转$2\pi/n$，$S$为关于一条固定对称轴的反射。变换
\[
T^k,\quad ST^k\qquad(0\leq k<n)
\]
都保持正多边形。$n$个旋转$T^k$的旋转角在模$2\pi$意义下不同，故彼此不同；左复合$S$可逆，故$n$个$ST^k$也彼此不同。旋转保持平面定向，反射与旋转的复合反转平面定向，所以两类不相交。已经得到$2n$个不同对称变换，结合前面的上界，得$|D_n|=2n$。

\textbf{再证明它不交换。}反射反转旋转方向，所以$STS=T^{-1}$，且$S^2=e$。若$ST=TS$，则右乘$S$得$STS=T$，与前式比较可得$T=T^{-1}$，即$T^2=e$。
但是$T$的阶恰为$n$：旋转$k$次恢复原位置当且仅当$2k\pi/n$为$2\pi$的整数倍，最小正整数$k$为$n$。所以$T^2=e$将给出$n\mid2$，与$n>2$矛盾。故$ST\ne TS$，$D_n$为所需的非交换群。
\end{proof}

\paragraph{复习提示}
本笔记约定$|D_n|=2n$。条件$n>2$不可略去，2阶和4阶群均交换。


\subsection{习题7：\texorpdfstring{反复共轭}{反复共轭}}
\label{ex:1:7}

\paragraph{题目}
若群$G$中的$a,b$满足$a^{-1}ba=b^r$，其中$r$为正整数，证明对任意正整数$i$都有$a^{-i}ba^i=b^{r^i}$。

\paragraph{考点与思路}
先验证共轭保持正整数次幂，再对共轭次数$i$归纳。

\paragraph{所用结论与原文依据}
\noteref{powers}；\noteref{glaws}。补充前置工具：数学归纳法。

% 原解来源：chapters/Part 2/chapter1-exercises.tex:58；第1章原\exref{1}{7}。

\begin{proof}
对任意正整数$m$，将相邻的$aa^{-1}$消去，得
\[
(a^{-1}ba)^m=a^{-1}b^ma.
\]
因此$a^{-1}b^ma=(b^r)^m=b^{rm}$。

当$i=1$时，结论就是题设。假设某个$i\geq1$已有$a^{-i}ba^i=b^{r^i}$，则
\begin{align*}
a^{-(i+1)}ba^{i+1}
&=a^{-1}(a^{-i}ba^i)a\\
&=a^{-1}b^{r^i}a\\
&=b^{r\cdot r^i}=b^{r^{i+1}}.
\end{align*}
归纳法给出所有正整数$i$的结论。
\end{proof}

\paragraph{复习提示}
反复共轭使指数每次乘以$r$，最终指数为$r^i$，不是$ri$。


\subsection{习题8：\texorpdfstring{取逆映射的同态性}{取逆映射的同态性}}
\label{ex:1:8}

\paragraph{题目}
证明：群$G$交换，当且仅当映射$\iota:G\to G$、$\iota(a)=a^{-1}$是群同态。取逆映射本身为双射，所以本题中的同态条件也等价于它是同构。

\paragraph{考点与思路}
取逆映射总为双射；真正需要判别的是它是否保乘法。

\paragraph{所用结论与原文依据}
\noteref{glaws}；\noteref{hom}。

% 原解来源：chapters/Part 2/chapter1-exercises.tex:68；第1章原\exref{1}{8}。

\begin{proof}
因为$\iota(\iota(a))=(a^{-1})^{-1}=a$，所以$\iota$的逆映射就是自身，因而总是双射。

\textbf{先由同态推出交换性。}若$\iota$为同态，则保乘法给出
\[
a^{-1}b^{-1}=\iota(a)\iota(b)=\iota(ab)=b^{-1}a^{-1}.
\]
两边取逆得到$ba=ab$。因$a,b$任意，$G$交换。

\textbf{再由交换性推出同构。}若$G$交换，则对任意$a,b\in G$，
\[
\iota(ab)=(ab)^{-1}=b^{-1}a^{-1}=a^{-1}b^{-1}=\iota(a)\iota(b).
\]
所以$\iota$为同态，再结合双射性，知它是同构。
\end{proof}

\paragraph{复习提示}
任何群都满足$(ab)^{-1}=b^{-1}a^{-1}$；保乘法所要求的却是$(ab)^{-1}=a^{-1}b^{-1}$。


\subsection{习题9：\texorpdfstring{由子集定义等价关系}{由子集定义等价关系}}
\label{ex:1:9}

\paragraph{题目}
设$S\subseteq G$非空，在$G$上定义$a\sim b\iff ab^{-1}\in S$。证明该关系为等价关系，当且仅当$S$为子群。

\paragraph{考点与思路}
等价关系要求自反、对称、传递。反向将关系与单位元联系起来，导出子群判别条件。

\paragraph{所用结论与原文依据}
\noteref{subgroup}；\noteref{cosets}。

% 原解来源：chapters/Part 2/chapter1-exercises.tex:76；第1章原\exref{1}{9}。

\begin{proof}
\textbf{若$S$为子群。}任取$a\in G$，因$aa^{-1}=e\in S$，有$a\sim a$，满足自反性。
若$a\sim b$，则$ab^{-1}\in S$，子群对逆封闭给出$ba^{-1}=(ab^{-1})^{-1}\in S$，所以$b\sim a$。
若$a\sim b$且$b\sim c$，则$ab^{-1},bc^{-1}\in S$，因而
\[
ac^{-1}=(ab^{-1})(bc^{-1})\in S,
\]
即$a\sim c$，满足传递性。因此该关系是等价关系。

\textbf{若该关系为等价关系。}自反性给出$e=aa^{-1}\in S$。任取$u,v\in S$，有$u\sim e$、$v\sim e$。由对称性得$e\sim v$，再由传递性得$u\sim v$，故$uv^{-1}\in S$。

说明此条件为何足够：取$u=e$可知$v^{-1}\in S$；再将已经得到的$v^{-1}$代入上述差式条件，得到$u(v^{-1})^{-1}=uv\in S$。所以$S$含单位元并对逆和乘积封闭，确为子群。

此外，$a\sim b$等价于$a=sb$对某个$s\in S$成立，因此$b$所在的等价类为右陪集$Sb$。
\end{proof}

\paragraph{复习提示}
$ab^{-1}$对应右陪集$Sb$；不要与$a^{-1}b$所对应的左陪集混淆。


\subsection{习题10：\texorpdfstring{$n\Z\cong\Z$}{nZ同构于Z}}
\label{ex:1:10}

\paragraph{题目}
设$n$为正整数，证明加法群$n\Z$与$\Z$同构。

\paragraph{考点与思路}
同构必须写出映射，并分别验证保运算、单射和满射。

\paragraph{所用结论与原文依据}
\noteref{iso}；\noteref{ringiso}。

% 原解来源：chapters/Part 2/chapter1-exercises.tex:88；第1章原\exref{1}{10}。

\begin{proof}
定义$\varphi:\Z\to n\Z$为$\varphi(k)=nk$。对任意$k,l\in\Z$，
\[
\varphi(k+l)=n(k+l)=nk+nl=\varphi(k)+\varphi(l),
\]
所以它保加法。若$\varphi(k)=\varphi(l)$，则$n(k-l)=0$。因为$n\ne0$，整数乘法无零因子，故$k=l$，所以它为单射。
任取$n\Z$的元素，按定义可写成$nk$，它就是$\varphi(k)$，所以映射满射。因此$\Z\cong n\Z$。
\end{proof}

\paragraph{复习提示}
这里讨论加法群。若$n>1$，上述映射不保通常的乘法，例如$\varphi(1\cdot1)=n$，而$\varphi(1)\varphi(1)=n^2$。若允许$n=0$，同构结论不成立。


\subsection{习题11：\texorpdfstring{克莱因四元群}{克莱因四元群}}
\label{ex:1:11}

\paragraph{题目}
证明$B=\{e,(12)(34),(13)(24),(14)(23)\}$为$S_4$的子群，并证明$B$不同构于四次单位根乘法群$U_4=\{1,i,-1,-i\}$。

\paragraph{考点与思路}
先用乘法关系验证子群，再用元素的阶排除与四次单位根群同构。

\paragraph{所用结论与原文依据}
\noteref{subgroup}；\noteref{iso}；\noteref{cycles}。

% 原解来源：chapters/Part 2/chapter1-exercises.tex:96；第1章原\exref{1}{11}。

\begin{proof}
令$u=(12)(34)$、$v=(13)(24)$、$w=(14)(23)$。分别作用于$1,2,3,4$可得
\[
u^2=v^2=w^2=e,\qquad uv=vu=w.
\]
由于$w=uv$且$u,v$交换，还可得到$uw=wu=v$、$vw=wv=u$。所以$B=\{e,u,v,w\}$对乘法封闭，每个元素的逆元仍属于$B$；由\proofref{subgroup}{子群判别法}，$B\leq S_4$。这就是克莱因四元群，记作$V_4$。

另一方面，四次单位根群$U_4=\{1,i,-1,-i\}$含有4阶元素$i$，而$B$中三个非单位元的阶均为2。同构保持元素的阶，因此$B\not\cong U_4$。
\end{proof}

\paragraph{复习提示}
两个群阶相同，只是同构的必要条件；还需比较元素阶或具体乘法。


\subsection{习题12：\texorpdfstring{指数2的子群正规}{指数2的子群正规}}
\label{ex:1:12}

\paragraph{题目}
若$G$为$2n$阶群，其中有一个$n$阶子群$H$，证明$H$正规。原题的条件给出$[G:H]=2$；以下按现有笔记的左右陪集方法证明这一更一般的指数2结论。

\paragraph{考点与思路}
指数2意味着除$H$以外，只有一个左陪集和一个右陪集，它们都必须是$H$的补集。

\paragraph{所用结论与原文依据}
\noteref{cosets}；\noteref{normal}。

% 原解来源：chapters/Part 2/chapter1-exercises.tex:106；第1章原\exref{1}{12}。

\begin{proof}
任取$a\in G$。若$a\in H$，子群性质给出$aH=H=Ha$。

若$a\notin H$，则$aH\ne H$，并且不同左陪集不交。因为恰有两个左陪集，必有$aH=G\setminus H$。
取逆映射将左陪集$xH$一一对应到右陪集$Hx^{-1}$，所以右陪集也恰有两个。同理$Ha=G\setminus H$。于是这一情形下也有$aH=Ha$。

对所有$a\in G$均已证明$aH=Ha$，故$H$正规。有限群的特殊情形由$[G:H]=|G|/|H|=2$立即得到。
\end{proof}

\paragraph{复习提示}
一般结论不要求群有限；只有用阶数比计算指数时才使用有限性。


\subsection{习题13：\texorpdfstring{偶数阶群有2阶元素}{偶数阶群有2阶元素}}
\label{ex:1:13}

\paragraph{题目}
证明偶数阶有限群中存在2阶元素。

\paragraph{考点与思路}
按取逆把元素配对，再利用群阶的奇偶性。

\paragraph{所用结论与原文依据}
\noteref{order}；\noteref{glaws}。

% 原解来源：chapters/Part 2/chapter1-exercises.tex:114；第1章原\exref{1}{13}。

\begin{proof}
设$|G|$为偶数。令$T=\{x\in G:x\ne x^{-1}\}$。若$x\in T$，则$x^{-1}\in T$；取逆将$T$分成互不相交的二元集合$\{x,x^{-1}\}$，所以$|T|$为偶数。

因此$G\setminus T=\{x\in G:x=x^{-1}\}$的元素个数也是偶数。该集合已含单位元$e$，故其元素个数是正偶数，至少为2。可取$a\ne e$使$a=a^{-1}$，于是$a^2=e$。

$a\ne e$排除了1阶情形，而$a^2=e$说明其阶整除2，故$\ord(a)=2$。
\end{proof}

\paragraph{复习提示}
取逆的不动点包括单位元，须由偶数计数证明还存在另一个不动点。


\subsection{习题14：\texorpdfstring{二面体群的矩阵实现}{二面体群的矩阵实现}}
\label{ex:1:14}

\paragraph{题目}
对$n\geq3$，令$\zeta=e^{2\pi i/n}$，
\[
A=\begin{pmatrix}0&1\\1&0\end{pmatrix},\qquad B=\begin{pmatrix}\zeta&0\\0&\zeta^{-1}\end{pmatrix}.
\]
则$\{B^k,AB^k:0\leq k<n\}\cong D_n$。

\paragraph{考点与思路}
核对矩阵满足二面体群的关系，并证明所列的两类矩阵恰有$2n$个。

\paragraph{所用结论与原文依据}
\noteref{dihedral}；\noteref{subgroup}；\noteref{iso}。

% 原解来源：chapters/Part 2/chapter1-exercises.tex:123；第1章原\exref{1}{14}。

\begin{proof}
直接进行矩阵乘法可得
\[
A^2=E,\qquad B^n=E,\qquad AB=B^{-1}A.
\]
$\zeta$的阶为$n$，故$B$的阶也为$n$。令$H=\{B^k,AB^k:0\leq k<n\}$。由$B^iA=AB^{-i}$，任意两个规定形状的矩阵相乘，分别有
\begin{align*}
B^iB^j&=B^{i+j},& B^i(AB^j)&=AB^{j-i},\\
(AB^i)B^j&=AB^{i+j},&(AB^i)(AB^j)&=B^{j-i}.
\end{align*}
指数均可模$n$化简，所以$H$对乘法封闭。又$(B^i)^{-1}=B^{-i}$，且$(AB^i)^2=E$，所以$H$对取逆封闭，且含$E$，故为群。

$B^0,\ldots,B^{n-1}$两两不同；左乘可逆矩阵$A$可知$AB^0,\ldots,AB^{n-1}$也两两不同。前一类是可逆对角矩阵，后一类是可逆反对角矩阵，两类不交，因而$|H|=2n$。

原笔记中$D_n$的元素为$T^k,ST^k$，并满足$T^n=S^2=e$及$ST=T^{-1}S$。定义$T^k\mapsto B^k$、$ST^k\mapsto AB^k$，所得映射为双射；两边乘法的四种情形均由上述相同关系给出，所以它保持乘法，故$D_n\cong H$。
\end{proof}

\paragraph{复习提示}
验证同构时，既要核对关系，也要说明正规形的唯一性和映射的双射性。


\subsection{习题15：\texorpdfstring{连续三个指数的乘积公式}{连续三个指数的乘积公式}}
\label{ex:1:15}

\paragraph{题目}
固定正整数$i$。若对所有$a,b\in G$及$k=i,i+1,i+2$均有$(ab)^k=a^kb^k$，则$G$交换。

\paragraph{考点与思路}
分别比较相邻的两个指数公式，得到$a$与$b^i,b^{i+1}$交换。

\paragraph{所用结论与原文依据}
\noteref{glaws}；\noteref{powers}。

% 原解来源：chapters/Part 2/chapter1-exercises.tex:137；第1章原\exref{1}{15}。

\begin{proof}
任取$a,b\in G$。比较指数$i$和$i+1$的公式，得到
\[
a^{i+1}b^{i+1}=(ab)^{i+1}=(ab)^i ab=a^ib^iab.
\]
在等式两边先左消去$a^i$，再右消去$b$，即得$ab^i=b^ia$。再把$i$换为$i+1$，比较指数$i+1,i+2$，同样得$ab^{i+1}=b^{i+1}a$。

由$a$与$b^i$交换，有$a(b^i)^{-1}=(b^i)^{-1}a$。于是
\begin{align*}
ab&=a(b^i)^{-1}b^{i+1}\\
&=(b^i)^{-1}ab^{i+1}\\
&=(b^i)^{-1}b^{i+1}a=ba.
\end{align*}
因为$a,b$任意，$G$为交换群。
\end{proof}

\paragraph{复习提示}
只比较两个相邻指数先得到与某个幂交换；第三个指数使这两个幂的指数差为1。


\subsection{习题16：\texorpdfstring{有限阶元素的积可以无限阶}{有限阶元素的积可以无限阶}}
\label{ex:1:16}

\paragraph{题目}
在$\SL_2(\Q)$中令
\[
A=\begin{pmatrix}0&-1\\1&0\end{pmatrix},\qquad B=\begin{pmatrix}0&1\\-1&-1\end{pmatrix}.
\]
则$A$为4阶，$B$为3阶，$AB$为无限阶。

\paragraph{考点与思路}
直接计算两个矩阵的幂，再把乘积的所有正整数次幂写出来。

\paragraph{所用结论与原文依据}
\noteref{order}。

% 原解来源：chapters/Part 2/chapter1-exercises.tex:147；第1章原\exref{1}{16}。

\begin{proof}
首先$\det A=\det B=1$，故$A,B\in\SL_2(\Q)$。矩阵乘法给出
\[
A^2=-E,\quad A^4=E,\qquad
B^2=\begin{pmatrix}-1&-1\\1&0\end{pmatrix},\quad B^3=E.
\]
$A\ne E$且$A^2\ne E$，所以$A$的阶恰为4；$B\ne E$且$B^3=E$，所以$B$的阶恰为3。

另一方面，
\[
AB=\begin{pmatrix}1&1\\0&1\end{pmatrix}.
\]
对$m$归纳：当$m=1$公式成立；若$(AB)^m=\left(\begin{smallmatrix}1&m\\0&1\end{smallmatrix}\right)$，则再乘$AB$得$(AB)^{m+1}=\left(\begin{smallmatrix}1&m+1\\0&1\end{smallmatrix}\right)$。因此对每个正整数$m$，
\[
(AB)^m=\begin{pmatrix}1&m\\0&1\end{pmatrix}\ne E.
\]
故$AB$为无限阶元素。
\end{proof}

\paragraph{复习提示}
群中两个有限阶元素的积未必有有限阶，交换性在相关结论中不能省略。


\subsection{习题17：\texorpdfstring{交换群的有限阶元素成子群}{交换群的有限阶元素成子群}}
\label{ex:1:17}

\paragraph{题目}
设$G$为交换群。证明$G$中全部有限阶元素组成一个子群，并说明交换性不能删去。

\paragraph{考点与思路}
使用子群判别法；有限阶元素之差的乘法记号是$ab^{-1}$。

\paragraph{所用结论与原文依据}
\noteref{subgroup}；\noteref{powers}；\exref{1}{16}。

% 原解来源：chapters/Part 2/chapter1-exercises.tex:162；第1章原\exref{1}{17}。

\begin{proof}
记$T$为$G$中全部有限阶元素的集合。单位元有1阶，故$T\ne\varnothing$。

任取$a,b\in T$，存在正整数$m,n$使$a^m=b^n=e$。由于$G$交换，$a$与$b^{-1}$也交换，故
\[
(ab^{-1})^{mn}=a^{mn}b^{-mn}=(a^m)^n(b^n)^{-m}=e.
\]
因此$ab^{-1}$的阶有限，属于$T$。由\proofref{subgroup}{子群判别法}，$T\leq G$。

交换性不能去掉：第一章习题16中$A,B$均为有限阶，但$AB$为无限阶。因此在该非交换群中，全部有限阶元素组成的集合对乘法不封闭。
\end{proof}

\paragraph{复习提示}
证明有限阶只需找到一个正幂等于单位元，不必计算其最小正指数。


\subsection{习题18：\texorpdfstring{只有有限多个子群的群必有限}{只有有限多个子群的群必有限}}
\label{ex:1:18}

\paragraph{题目}
证明：如果群$G$只有有限多个子群，那么$G$必为有限群。

\paragraph{考点与思路}
先排除无限阶元素，再将全群表示为有限多个有限循环子群的并。

\paragraph{所用结论与原文依据}
\noteref{cyclic}。

% 原解来源：chapters/Part 2/chapter1-exercises.tex:170；第1章原\exref{1}{18}。

\begin{proof}
\textbf{第一步：每个元素都为有限阶。}
反设$a\in G$为无限阶，则子群$\langle a^m\rangle$（$m=1,2,\ldots$）两两不同。
事实上，若$\langle a^m\rangle=\langle a^n\rangle$，则存在整数$u,v$使$a^m=a^{nu}$、$a^n=a^{mv}$。无限阶意味着相等幂的指数相等，故$m=nu$、$n=mv$。由于$m,n$为正整数，$u,v$均为正整数；于是$uv=1$，得$u=v=1$、$m=n$。
这产生无限多个子群，与假设矛盾。

\textbf{第二步：全群有限。}
所有循环子群均为$G$的子群，所以它们只有有限多个，记为$C_1,\ldots,C_t$。第一步说明每个$C_j$都有限。
任意$g\in G$属于其生成的循环子群$\langle g\rangle$，因此
\[
G=\bigcup_{j=1}^{t}C_j,\qquad |G|\leq\sum_{j=1}^{t}|C_j|<\infty.
\]
所以$G$有限。
\end{proof}

\paragraph{复习提示}
仅证明“每个元素有限阶”还不足以证明群有限；必须利用循环子群的个数也有限。


\subsection{习题19：\texorpdfstring{$D_n$的全部正规子群}{Dn的全部正规子群}}
\label{ex:1:19}

\paragraph{题目}
使用正文中$|D_n|=2n$的约定，$n\geq3$。

\paragraph{考点与思路}
先分清只含旋转的子群与含反射的子群，再用正规性限制参数。

\paragraph{所用结论与原文依据}
\noteref{dihedral}；\noteref{cyclic}；\noteref{normal}；\exref{1}{12}。

% 原解来源：chapters/Part 2/chapter1-exercises.tex:179；第1章原\exref{1}{19}。

\begin{proof}
沿用$D_n=\{T^k,ST^k:0\leq k<n\}$及$ST=T^{-1}S$。设$H\leq D_n$。

先说明旋转子群中的子群形式。设$K\leq\langle T\rangle$，因为$T^n=e\in K$，可取最小正整数$d$使$T^d\in K$。若$T^l\in K$，把$l$写成$qd+r$，$0\leq r<d$，则$T^r=T^l(T^d)^{-q}\in K$。由$d$的最小性只能$r=0$，所以$K=\langle T^d\rangle$；对$l=n$使用同一结论得$d\mid n$。

若$H\subseteq\langle T\rangle$，由上段可写$H=\langle T^d\rangle$，$d\mid n$。用$T$共轭不改变该子群，用$S$共轭把$T^d$变为$T^{-d}$，仍生成同一子群。由于$S,T$生成$D_n$，这样的子群全都正规。

若$H$含反射$ST^k$，写$H\cap\langle T\rangle=\langle T^d\rangle$，$d\mid n$。对$H$中任意另一反射$ST^l$，有
\[
(ST^k)(ST^l)=T^{l-k}\in H\cap\langle T\rangle,
\]
所以$d\mid l-k$。因此
\[
H=\langle T^d,ST^k\rangle,
\]
且$k$只需模$d$取值。这也证明上述形式已穷尽含反射的子群。

若这样的$H$正规，则$T(ST^k)T^{-1}=ST^{k-2}\in H$，由刚才的判据得$d\mid2$。若$d=1$，则$H=D_n$；若$d=2$，须$n$为偶数，且$k$模2只有0、1两种，得到
\[
\langle T^2,S\rangle,\qquad\langle T^2,ST\rangle.
\]
这两个子群均有$n$个元素，指数为2，由第一章习题12可知它们正规。

综上，全部正规子群为$\langle T^d\rangle$（$d\mid n$）、$D_n$，以及$n$为偶数时另有的上述两个指数2子群。
\end{proof}

\paragraph{复习提示}
必要条件$d\mid2$得到候选后，还须验证每个候选确实正规。


\subsection{习题20：\texorpdfstring{子群乘积的判别}{子群乘积的判别}}
\label{ex:1:20}

\paragraph{题目}
对$H,K\leq G$，$HK$为子群当且仅当$HK=KH$。

\paragraph{考点与思路}
用集合取逆公式证明必要性，再分别验证乘法和取逆封闭。

\paragraph{所用结论与原文依据}
\noteref{setproduct}；\noteref{subgroup}。

% 原解来源：chapters/Part 2/chapter1-exercises.tex:199；第1章原\exref{1}{20}。

\begin{proof}
若$HK$为子群，则它对取逆封闭，因而
\[
HK=(HK)^{-1}=K^{-1}H^{-1}=KH.
\]
这里使用$H^{-1}=H$、$K^{-1}=K$。

反之，设$HK=KH$。$e=ee\in HK$，故$HK$非空。集合乘法的结合律给出
\begin{align*}
(HK)(HK)&=H(KH)K=H(HK)K\\
&=(HH)(KK)=HK.
\end{align*}
所以$HK$对乘法封闭。又$(HK)^{-1}=KH=HK$，故对取逆封闭。由\proofref{subgroup}{子群判别法}，$HK$为子群。
\end{proof}

\paragraph{复习提示}
集合相等$HK=KH$不要求每个$h\in H$与每个$k\in K$交换。


\subsection{习题21：\texorpdfstring{有限子群乘积的计数}{有限子群乘积的计数}}
\label{ex:1:21}

\paragraph{题目}
设$H,K$为群$G$的有限子群，证明
\[
|HK|=\frac{|H||K|}{|H\cap K|}.
\]
此处不要求集合$HK$是子群。

\paragraph{考点与思路}
映射$(h,k)\mapsto hk$一般不单射；准确计算每个像的原像个数。

\paragraph{所用结论与原文依据}
\noteref{cosets}。补充前置工具：有限满射的等纤维计数（证明中给出纤维双射）。

% 原解来源：chapters/Part 2/chapter1-exercises.tex:208；第1章原\exref{1}{21}。

\begin{proof}
定义满射$\pi:H\times K\to HK$，$\pi(h,k)=hk$。固定$g=hk\in HK$的一种表示。

若$h'k'=hk$，两边左乘$h^{-1}$、右乘$k'^{-1}$，得到
\[
t:=h^{-1}h'=kk'^{-1}\in H\cap K.
\]
由此$h'=ht$、$k'=t^{-1}k$。反之，任取$t\in H\cap K$，都有$ht\in H$、$t^{-1}k\in K$，且
\[
(ht)(t^{-1}k)=hk=g.
\]
不同$t$给出不同的有序对，因为$ht_1=ht_2$可消去$h$而得到$t_1=t_2$。因此$g$的全部原像与$H\cap K$一一对应，恰有$|H\cap K|$个。

对有限集合$H\times K$按原像分组计数，得
\[
|H||K|=|HK|\,|H\cap K|,
\]
除以非零正整数$|H\cap K|$即得所求。
\end{proof}

\paragraph{复习提示}
需要扣除的是交集造成的重复表示；即使$HK$不是子群，集合计数公式仍成立。


\subsection{习题22：\texorpdfstring{正规子群的乘积}{正规子群的乘积}}
\label{ex:1:22}

\paragraph{题目}
设$M,N\trianglelefteq G$，证明$MN=NM$且$MN\trianglelefteq G$。若$M\cap N=\{e\}$，证明$MN/N\cong M$。

\paragraph{考点与思路}
先证明集合乘积是正规子群，再构造商群的满同态并求核。

\paragraph{所用结论与原文依据}
\noteref{normal}；\noteref{firstiso}；\noteref{setproduct}；\exref{1}{20}。

% 原解来源：chapters/Part 2/chapter1-exercises.tex:221；第1章原\exref{1}{22}。

\begin{proof}
\textbf{第一步：证明$MN=NM$。}
任取$m\in M,n\in N$，因$M$正规，$n^{-1}mn\in M$，于是$mn=n(n^{-1}mn)\in NM$，故$MN\subseteq NM$。
同样，由$M$正规，$nmn^{-1}\in M$，所以$nm=(nmn^{-1})n\in MN$，得到反向包含。

\textbf{第二步：证明$MN$为正规子群。}
$e\in MN$。集合乘积满足结合律，利用$NM=MN$，得
\[
(MN)(MN)=M(NM)N=M(MN)N=MN,
\qquad (MN)^{-1}=NM=MN.
\]
所以$MN$对乘法和逆封闭，为子群。再任取$g\in G$，由两个子群的正规性，
\[
g(MN)g^{-1}=(gMg^{-1})(gNg^{-1})=MN.
\]
因此$MN\trianglelefteq G$。

\textbf{第三步：构造同构。}
因为$N$在$G$中正规，它在子群$MN$中也正规，可以构成商群$MN/N$。定义
\[
\varphi:M\longrightarrow MN/N,\qquad \varphi(m)=mN.
\]
商群运算给出$\varphi(m_1m_2)=m_1m_2N=(m_1N)(m_2N)$，所以$\varphi$是同态。
任意陪集$mnN$等于$mN$，故$\varphi$满射。其核为
\[
\ker\varphi=\{m\in M:mN=N\}=M\cap N.
\]
因此\proofref{firstiso}{同态基本定理}给出$M/(M\cap N)\cong MN/N$。当$M\cap N=\{e\}$时，$\varphi$本身单射，故$M\cong MN/N$。
\end{proof}

\paragraph{复习提示}
证明商群同构时，应逐项交代映射、保运算、满射和核；交平凡正是单射的依据。


\subsection{习题23：\texorpdfstring{中心化子与正规化子}{中心化子与正规化子}}
\label{ex:1:23}

\paragraph{题目}
对非空$S\subseteq G$，定义
\[
C_G(S)=\{x:xa=ax\text{对所有 }a\in S\},\quad
N_G(S)=\{x:x^{-1}Sx=S\}.
\]
则两者均为子群，且$C_G(S)\trianglelefteq N_G(S)$。

\paragraph{考点与思路}
先分别验证两个子群，再将中心化子中的元素作共轭。

\paragraph{所用结论与原文依据}
\noteref{subgroup}；\noteref{normal}。

% 原解来源：chapters/Part 2/chapter1-exercises.tex:232；第1章原\exref{1}{23}。

\begin{proof}
两集合均含$e$。若$x,y\in C_G(S)$，则对任意$a\in S$，$y^{-1}a=ay^{-1}$，从而
\[
(xy^{-1})a=xay^{-1}=axy^{-1}.
\]
故$xy^{-1}\in C_G(S)$，\proofref{subgroup}{子群判别法}给出$C_G(S)\leq G$。

若$x,y\in N_G(S)$，则$x^{-1}Sx=S$、$y^{-1}Sy=S$，后式也给出$ySy^{-1}=S$。于是
\[
(xy^{-1})^{-1}S(xy^{-1})=y(x^{-1}Sx)y^{-1}=ySy^{-1}=S.
\]
所以$N_G(S)$也是子群。若$c\in C_G(S)$，则对所有$a\in S$有$c^{-1}ac=a$，故$c^{-1}Sc=S$，说明$C_G(S)\leq N_G(S)$。

最后任取$n\in N_G(S)$、$c\in C_G(S)$。因为$n^{-1}an\in S$，对每个$a\in S$都有
\begin{align*}
(ncn^{-1})a&=nc(n^{-1}an)n^{-1}\\
&=n(n^{-1}an)cn^{-1}=a(ncn^{-1}).
\end{align*}
故$ncn^{-1}\in C_G(S)$。由\proofref{normal}{正规性判别}，$C_G(S)\trianglelefteq N_G(S)$。
\end{proof}

\paragraph{复习提示}
正规化子只要求共轭后整个集合不变；中心化子要求其中每个元素都被固定。


\subsection{习题24：\texorpdfstring{2阶群}{2阶群}}
\label{ex:1:24}

\paragraph{题目}
证明任意2阶群均同构于乘法群$\{1,-1\}$。

\paragraph{考点与思路}
用拉格朗日定理确定非单位元的阶，直接写出同构。

\paragraph{所用结论与原文依据}
\noteref{lagrange}；\noteref{iso}。

% 原解来源：chapters/Part 2/chapter1-exercises.tex:252；第1章原\exref{1}{24}。

\begin{proof}
设$G=\{e,a\}$。由\proofref{lagrange}{拉格朗日定理}，$\ord(a)\mid2$；$a\ne e$，所以$\ord(a)=2$，即$a^2=e$。

定义$\varphi:G\to\{1,-1\}$为$\varphi(e)=1$、$\varphi(a)=-1$。它是双射。涉及$e$的乘积显然保持；剩下的$a^2=e$对应$(-1)^2=1$，所以$\varphi$也保持这一乘积，故为群同构。
\end{proof}

\paragraph{复习提示}
两元素群的乘法由单位元性质和$a^2=e$唯一确定。


\subsection{习题25：\texorpdfstring{4阶群的分类}{4阶群的分类}}
\label{ex:1:25}

\paragraph{题目}
确定4阶群的全部同构类型。

\paragraph{考点与思路}
按是否存在4阶元素分类；不存在时确定全部乘法。

\paragraph{所用结论与原文依据}
\noteref{lagrange}；\noteref{cyclic}；\noteref{iso}；\exref{1}{2}。

% 原解来源：chapters/Part 2/chapter1-exercises.tex:252；第1章原\exref{1}{25}。

\begin{proof}
设$|G|=4$。若存在4阶元素$a$，则$\langle a\rangle$已有4个元素，所以$G=\langle a\rangle\cong\Z/4\Z$。

若不存在4阶元素，由\proofref{lagrange}{拉格朗日定理}，每个非单位元都只能为2阶，故所有$x\in G$均满足$x^2=e$。第一章习题2给出$G$交换。取两个不同的非单位元$a,b$。若$ab=e$，则$b=a^{-1}=a$，矛盾；若$ab=a$或$ab=b$，消去得$b=e$或$a=e$，也矛盾。因此$e,a,b,ab$两两不同，已穷尽$G$。

由$a^2=b^2=e$及$ab=ba$，映射
\[
(\bar i,\bar j)\longmapsto a^ib^j\quad
\text{从 }(\Z/2\Z)\times(\Z/2\Z)\text{ 到 }G
\]
良定义、保乘法且为双射，故$G\cong V_4$。

$\Z/4\Z$含4阶元素，而$V_4$没有，两者不同构。这证明4阶群恰有上述两种同构类型。
\end{proof}

\paragraph{复习提示}
分类必须穷尽，并且最后要证明两个候选不同构。


\subsection{习题26：\texorpdfstring{素数阶群}{素数阶群}}
\label{ex:1:26}

\paragraph{题目}
证明任意两个相同素数阶的群同构。

\paragraph{考点与思路}
素数阶使每个非单位元都生成全群，再写出循环群的同构。

\paragraph{所用结论与原文依据}
\noteref{lagrange}；\noteref{cyclic}；\noteref{iso}。

% 原解来源：chapters/Part 2/chapter1-exercises.tex:252；第1章原\exref{1}{26}。

\begin{proof}
设$|G|=p$，其中$p$为素数。任取$a\ne e$，由\proofref{lagrange}{拉格朗日定理}有$\ord(a)\mid p$，且$\ord(a)\ne1$，所以$\ord(a)=p$。因此$G=\langle a\rangle$。

定义$\varphi:\Z/p\Z\to G$，$\varphi(\bar k)=a^k$。若$k\equiv l\pmod p$，则$a^{k-l}=e$，所以$a^k=a^l$，映射良定义。又
\[
\varphi(\bar k+\bar l)=a^{k+l}=a^ka^l,
\]
故为同态。它满射；若$a^k=e$，则$p\mid k$，故核为$\{\bar0\}$，所以也单射。因此$G\cong\Z/p\Z$，任意两个$p$阶群均同构。
\end{proof}

\paragraph{复习提示}
素数阶结论中的素数$p$须相同，才可断言两个群同构。


\subsection{习题27：\texorpdfstring{整数对上的环}{整数对上的环}}
\label{ex:1:27}

\paragraph{题目}
在$\Z\times\Z$上定义
\[
(a,b)+(c,d)=(a+c,b+d),\qquad(a,b)(c,d)=(ac+bd,ad+bc).
\]

\paragraph{考点与思路}
沿原笔记把整数对嵌入矩阵环，借子环结构验证环公理。

\paragraph{所用结论与原文依据}
\noteref{subring}；\noteref{ringiso}；\noteref{zerodiv}。

% 原解来源：chapters/Part 2/chapter1-exercises.tex:269；第1章原\exref{1}{27}。

\begin{proof}
令
\[
\varphi(a,b)=\begin{pmatrix}a&b\\b&a\end{pmatrix},\qquad
S=\left\{\begin{pmatrix}a&b\\b&a\end{pmatrix}:a,b\in\Z\right\}.
\]
$S$含零矩阵，两个此类矩阵之差仍有同样形状，而
\[
\begin{pmatrix}a&b\\b&a\end{pmatrix}
\begin{pmatrix}c&d\\d&c\end{pmatrix}
=\begin{pmatrix}ac+bd&ad+bc\\ad+bc&ac+bd\end{pmatrix}\in S.
\]
故由\proofref{subring}{子环判别法}，$S\leq M_2(\Z)$。映射$\varphi$是$\Z\times\Z$到$S$的双射，并且按题设两种运算有
\[
\varphi((a,b)+(c,d))=\varphi(a,b)+\varphi(c,d),\qquad
\varphi((a,b)(c,d))=\varphi(a,b)\varphi(c,d).
\]
所以题设运算使$\Z\times\Z$成为与$S$同构的环。

由乘法公式交换两对参数不改变结果，环是交换的；加法零元为$(0,0)$，负元为$(-a,-b)$，乘法单位元为$(1,0)$。但
\[
(1,1)(1,-1)=(0,0),
\]
两个因子均非零，故此环有零因子，不是整环。
\end{proof}

\paragraph{复习提示}
使用矩阵模型前，必须确认像确实为子环，并且模型保持题设的两种运算。


\subsection{习题28：\texorpdfstring{交换整数加法与乘法的角色}{交换整数加法与乘法的角色}}
\label{ex:1:28}

\paragraph{题目}
令$a\oplus b=ab$，$a\odot b=a+b$，这是否使$\Z$成环？

\paragraph{考点与思路}
先检查新加法能否构成交换群，任何一个公理失败即可否定。

\paragraph{所用结论与原文依据}
\noteref{ring}；\noteref{group}。

% 原解来源：chapters/Part 2/chapter1-exercises.tex:281；第1章原\exref{1}{28}。

\begin{proof}
若$\Z$在新运算下为环，则$(\Z,\oplus)$必须为交换群。由于$a\oplus b=ab$，新加法的单位元必须是整数$1$：它满足$a\oplus1=1\oplus a=a$。

然而整数$0$没有新加法逆元，因为对任意$x\in\Z$都有$0\oplus x=0x=0\ne1$。因此$(\Z,\oplus)$不是群，题设两种运算不能使$\Z$成环。
\end{proof}

\paragraph{复习提示}
环的零元指加法单位元，换运算后应重新确定，不能沿用整数0。


\subsection{习题29：\texorpdfstring{改变坐标得到同构环}{改变坐标得到同构环}}
\label{ex:1:29}

\paragraph{题目}
在交换幺环$R$的底集上定义
\[
a\oplus b=a+b-1,\qquad a\odot b=a+b-ab.
\]
证明新运算使它成为与原环同构的交换幺环，并求新的零元、单位元及加法逆元。

\paragraph{考点与思路}
寻找坐标变换$f(a)=1-a$，把新运算转回熟悉的运算。

\paragraph{所用结论与原文依据}
\noteref{ringiso}；\noteref{ring}。

% 原解来源：chapters/Part 2/chapter1-exercises.tex:289；第1章原\exref{1}{29}。

\begin{proof}
记具有新运算的集合为$R_{\mathrm{new}}$，定义$f:R_{\mathrm{new}}\to R$为$f(a)=1-a$。因为$f(f(a))=1-(1-a)=a$，它是双射。
对任意$a,b\in R$，
\begin{align*}
f(a\oplus b)&=1-(a+b-1)=2-a-b=f(a)+f(b),\\
f(a\odot b)&=1-(a+b-ab)=1-a-b+ab\\
&=(1-a)(1-b)=f(a)f(b).
\end{align*}
所以$f$将两种新运算分别对应到原环的加法、乘法。由于$f$单射，原环的每一条运算恒等式均可通过$f$拉回；例如
\begin{align*}
f((a\odot b)\odot c)&=(f(a)f(b))f(c)\\
&=f(a)(f(b)f(c))=f(a\odot(b\odot c)),
\end{align*}
从而新乘法满足结合律。加法的结合、交换以及分配律同理由原环的相应公理得到。因此新结构为环，且$f$为环同构。

新零元是$f^{-1}(0)=1$，新乘法单位元是$f^{-1}(1)=0$。也可直接核对$a\oplus1=a$、$a\odot0=a$。元素$a$的新加法逆元为$2-a$，因为
\[
a\oplus(2-a)=a+(2-a)-1=1,
\]
右端正是新零元。新乘法交换性由$a+b-ab=b+a-ba$得到。
\end{proof}

\paragraph{复习提示}
式中$0,1,2$表示原环元素，$2=1+1$；新结构的零元与单位元分别是原来的$1$与$0$。


\subsection{习题30：\texorpdfstring{子环与母环的单位元和交换性}{子环与母环的单位元和交换性}}
\label{ex:1:30}

\paragraph{题目}
分别举例说明：(i)母环有单位元而子环无单位元；(ii)母环无单位元而子环有单位元；(iii)母环和子环都有单位元，但二者不同；(iv)母环不交换而子环交换。

\paragraph{考点与思路}
这里子环只要求对加法、取负和乘法封闭，并不要求与母环共用单位元。应检查实例确实满足对应性质。

\paragraph{所用结论与原文依据}
\noteref{subring}；\noteref{unit}。

% 原解来源：chapters/Part 2/chapter1-exercises.tex:302；第1章原\exref{1}{30}。

\begin{proof}
\begin{enumerate}\item 取$R=\Z$、$S=2\Z$。偶数对差和积封闭，所以$S$是子环；$R$的单位元为1。若$u\in2\Z$是$S$的单位元，则$u\cdot2=2$，在整数中得$u=1$，与$u$为偶数矛盾。因此$S$没有单位元。

\item 取逐分量运算的环$R=2\Z\times\Z$，以及子环$S=\{0\}\times\Z$。若$R$有单位元$(u,v)$，将它乘以$(2,0)$得到$(2u,0)=(2,0)$，所以$u=1\notin2\Z$，矛盾。而对任意$(0,k)\in S$，
\[
(0,1)(0,k)=(0,k)=(0,k)(0,1),
\]
所以$S$有单位元$(0,1)$。

\item 取$R=M_2(\R)$，令
\[
S=\left\{\begin{pmatrix}a&0\\0&0\end{pmatrix}:a\in\R\right\}.
\]
两个这样的矩阵之差和积仍有相同形状，所以$S$是子环。母环的单位元为$E=\operatorname{diag}(1,1)$，子环的单位元为$\operatorname{diag}(1,0)$，两者不同。

\item 仍取$R=M_2(\R)$，令$S=\{aE:a\in\R\}$为数量矩阵构成的子环。因$(aE)(bE)=abE=baE$，$S$交换。记$E_{ij}$为第$(i,j)$项为1、其余项为0的矩阵，则
\[
E_{12}E_{21}=E_{11},\qquad E_{21}E_{12}=E_{22},
\]
二者不同，故母环不交换。
\end{enumerate}
\end{proof}

\paragraph{复习提示}
单位元是相对于环内全部元素定义的；子环中的单位元不必对母环其他元素起单位作用。


\subsection{习题31：\texorpdfstring{单侧单位元}{单侧单位元}}
\label{ex:1:31}

\paragraph{题目}
设$e_La=a$对所有$a\in R$成立时称$e_L$为左单位元，右单位元类似。证明：
\begin{enumerate}
\item 若环同时有左单位元和右单位元，则两者相等，为双侧单位元。
\item 无零因子的环若有左单位元，则它也为右单位元。
\item 若环有左单位元而无右单位元，则左单位元不唯一。
\end{enumerate}

\paragraph{考点与思路}
分别验证三项结论，严格区别左单位元、右单位元和双侧单位元。

\paragraph{所用结论与原文依据}
\noteref{unit}；\noteref{zerodiv}；\noteref{ringlaws}。

% 原解来源：chapters/Part 2/chapter1-exercises.tex:315；第1章原\exref{1}{31}。

\begin{proof}
\begin{enumerate}
\item 设$e$为左单位元、$f$为右单位元。由$e$的性质得$ef=f$，由$f$的性质得$ef=e$，所以$e=f$。这个共同元素同时满足两侧单位元条件。

\item 设$R$无零因子且$e$为左单位元。若$R=\{0\}$，唯一元素就是双侧单位元。若$R\ne\{0\}$，固定$b\ne0$。对任意$a\in R$，
\[
(ae-a)b=a(eb)-ab=ab-ab=0.
\]
因为$b\ne0$且无零因子，必有$ae-a=0$。所以$ae=a$对每个$a$成立，$e$也为右单位元。

\item 设$e$为左单位元，而$R$没有右单位元。若所有$a$均满足$ae=a$，则$e$就是右单位元，矛盾。因此存在$a$使$d=a-ae\ne0$。对任意$x\in R$，
\[
dx=(a-ae)x=ax-a(ex)=0.
\]
于是$(e+d)x=ex+dx=x$。所以$e+d$也是左单位元，且因$d\ne0$有$e+d\ne e$，从而左单位元至少有两个。
\end{enumerate}
\end{proof}

\paragraph{复习提示}
无零因子的消去必须有一个非零因子；第二项要单独交代零环。


\subsection{习题32：\texorpdfstring{域没有非平凡双边理想}{域没有非平凡双边理想}}
\label{ex:1:32}

\paragraph{题目}
证明域$F$的双边理想只有$\{0\}$和$F$。同一结论也适用于体。

\paragraph{考点与思路}
非零理想中取非零元素，用其逆元推出单位元属于该理想。

\paragraph{所用结论与原文依据}
\noteref{ideal}；\noteref{field}；\noteref{idealunit}。

% 原解来源：chapters/Part 2/chapter1-exercises.tex:331；第1章原\exref{1}{32}。

\begin{proof}
设$I$为$F$的理想。若$I=\{0\}$，已经属于所列情形。若$I\ne\{0\}$，可取$a\in I$、$a\ne0$。域中$a$有逆元$a^{-1}\in F$，理想的吸收性给出$1=a^{-1}a\in I$。
任取$r\in F$，再次由吸收性，$r=r\cdot1\in I$，所以$F\subseteq I$。结合$I\subseteq F$，得到$I=F$。

证明只使用非零元素有双边逆元，没有使用乘法交换性，所以对不一定交换的体也成立。
\end{proof}

\paragraph{复习提示}
一般环中非零元素未必可逆，不能对任意环照搬这一步。


\subsection{习题33：\texorpdfstring{$Ra$何时是理想}{Ra何时是理想}}
\label{ex:1:33}

\paragraph{题目}
若$R$交换，则$Ra=\{ra:r\in R\}$是双边理想；非交换时未必。

\paragraph{考点与思路}
逐项验证理想条件；非交换反例仍取矩阵单位。

\paragraph{所用结论与原文依据}
\noteref{ideal}；\noteref{ringlaws}。

% 原解来源：chapters/Part 2/chapter1-exercises.tex:339；第1章原\exref{1}{33}。

\begin{proof}
设$R$交换，$a\in R$固定。$0=0a\in Ra$，且对$r,s\in R$，
\[
ra-sa=(r-s)a\in Ra.
\]
因此$Ra$为加法子群。对任意$t\in R$，结合律及交换律给出
\[
t(ra)=(tr)a\in Ra,\qquad (ra)t=(rt)a\in Ra.
\]
所以$Ra$是双边理想。

在非交换环$R=M_2(\Q)$中取$a=E_{11}$，其中$E_{ij}$表示仅第$(i,j)$个位置为1、其余为0的矩阵。直接计算$XE_{11}$可知，$Ra$恰为第二列为零的矩阵集合。$E_{11}\in Ra$，但
\[
E_{11}E_{12}=E_{12}\notin Ra.
\]
所以$Ra$不一定是双边理想。

若$R$没有单位元，不能直接断言$a\in Ra$；本题证明理想条件没有使用这一断言。
\end{proof}

\paragraph{复习提示}
$Ra$总有左吸收性；要成为双边理想，还要检查右吸收性。


\subsection{习题34：\texorpdfstring{理想的根}{理想的根}}
\label{ex:1:34}

\paragraph{题目}
设$R$为交换环、$I$为其理想，定义
\[
\sqrt I=\{r\in R:\text{存在正整数 }n\text{ 使 }r^n\in I\}.
\]
证明$\sqrt I$为$R$的理想。

\paragraph{考点与思路}
验证差封闭和吸收性。两个元素对应的幂次数可能不同，统一选取$n+m-1$次幂。

\paragraph{所用结论与原文依据}
\noteref{ideal}；\noteref{ringlaws}。补充前置工具：交换环的二项式展开；无幺环中的端项不使用零次幂。

% 原解来源：chapters/Part 2/chapter1-exercises.tex:349；第1章原\exref{1}{34}。

\begin{proof}
首先$0^1=0\in I$，故$0\in\sqrt I$，集合非空。

任取$r,s\in\sqrt I$，由定义可选正整数$n,m$使$r^n\in I$、$s^m\in I$。令$N=n+m-1$。若环无单位元，下面求和的两个端项分别指$(-s)^N$与$r^N$，不使用零次幂。因为环交换，可以使用二项式展开：
\[
(r-s)^N=\sum_{k=0}^{N}(-1)^{N-k}\binom Nk r^k s^{N-k}.
\]
对任一项，若$k\geq n$，它含有因子$r^n$，故属于$I$；若$k<n$，则$k\leq n-1$，从而$N-k\geq m$，这一项含有因子$s^m$，也属于$I$。整数倍与有限和仍在加法子群$I$中，故$(r-s)^N\in I$，即$r-s\in\sqrt I$。因此$\sqrt I$为加法子群。

任取$t\in R$、$r\in\sqrt I$，选$n\geq1$使$r^n\in I$。由交换性和理想吸收性，
\[
(tr)^n=t^nr^n\in I.
\]
所以$tr\in\sqrt I$；交换性又给出$rt=tr\in\sqrt I$。因此$\sqrt I$满足双边吸收性，确为理想。

若$R$没有单位元，上述展开仍然成立：在$k=0$或$k=N$的端项中，直接写作$(-s)^N$或$r^N$，不引入$R$内不存在的零次幂。判断因子属于$I$时，若剩余因子个数为零，就直接使用$r^n\in I$或$s^m\in I$。
\end{proof}

\paragraph{复习提示}
不能仅由$r^n,s^m\in I$断言$(r-s)^{\max(n,m)}\in I$；$n+m-1$保证展开中每一项都含足够多的某一种因子。


\subsection{习题35：\texorpdfstring{交换幺环仅有平凡理想时为域}{交换幺环仅有平凡理想时为域}}
\label{ex:1:35}

\paragraph{题目}
假设$1\neq0$。

\paragraph{考点与思路}
非零元素生成的主理想只能等于全环，因而可以解出逆元。

\paragraph{所用结论与原文依据}
\noteref{ideal}；\noteref{field}；\exref{1}{33}。

% 原解来源：chapters/Part 2/chapter1-exercises.tex:360；第1章原\exref{1}{35}。

\begin{proof}
设$R$为交换幺环，$1\ne0$，且只有$\{0\}$与$R$两个理想。任取$a\ne0$。由第一章习题33，$Ra$为理想；它含$a=1a\ne0$，所以$Ra\ne\{0\}$，只能$Ra=R$。

因此$1\in Ra$，存在$r\in R$使$ra=1$。由交换性也有$ar=1$，故$a$可逆。每个非零元素均可逆，按域的定义，$R$为域。

条件$1\ne0$不能略去：零环也只有平凡理想，但不是域。
\end{proof}

\paragraph{复习提示}
这里需要幺元来保证$a\in Ra$，需要交换性来把单侧逆变为双侧逆。


\subsection{习题36：\texorpdfstring{全矩阵环是单环}{全矩阵环是单环}}
\label{ex:1:36}

\paragraph{题目}
$M_n(\Q)$没有非平凡双边理想。

\paragraph{考点与思路}
从理想中一个非零矩阵的非零系数出发，用左右矩阵单位取得所有矩阵单位。

\paragraph{所用结论与原文依据}
\noteref{ideal}。

% 原解来源：chapters/Part 2/chapter1-exercises.tex:369；第1章原\exref{1}{36}。

\begin{proof}
设$I$为$M_n(\Q)$的非零双边理想，取$0\ne A=(a_{ij})\in I$。存在$i,j$使$a_{ij}\ne0$。记$E_{kl}$为矩阵单位，满足$E_{ij}E_{kl}=\delta_{jk}E_{il}$。

任取$k,l$，理想的两侧吸收性给出
\[
E_{ki}AE_{jl}=a_{ij}E_{kl}\in I.
\]
再左乘矩阵$a_{ij}^{-1}E_{kk}$，得$E_{kl}\in I$。因此所有矩阵单位都属于$I$。

对任意有理矩阵$X=(x_{kl})$，有$x_{kl}E_{kl}=(x_{kl}E_{kk})E_{kl}\in I$。理想对有限和封闭，所以
\[
X=\sum_{k,l}x_{kl}E_{kl}\in I.
\]
因此$I=M_n(\Q)$，不存在非零真双边理想。
\end{proof}

\paragraph{复习提示}
不能直接把理想当作有理向量子空间；数乘封闭应由母环矩阵的吸收性推出。


\subsection{习题37：\texorpdfstring{$aba=0$判别零因子}{aba=0判别零因子}}
\label{ex:1:37}

\paragraph{题目}
若$a,b\neq0$且$aba=0$，则$a$为左零因子或右零因子。

\paragraph{考点与思路}
按$ba$是否为零分类，保证出现的另一个零乘积因子非零。

\paragraph{所用结论与原文依据}
\noteref{zerodiv}。

% 原解来源：chapters/Part 2/chapter1-exercises.tex:380；第1章原\exref{1}{37}。

\begin{proof}
若$ba=0$，因为$b\ne0$、$a\ne0$，按右零因子的定义，$a$是右零因子。

若$ba\ne0$，则由结合律及题设，$a(ba)=aba=0$。此时$a$与$ba$均非零，按左零因子的定义，$a$是左零因子。

上述两种情形穷尽所有可能，故结论成立。
\end{proof}

\paragraph{复习提示}
本笔记称$ab=0$中的$a$为左零因子；引用时沿用这一左右约定。


\subsection{习题38：\texorpdfstring{$1-a$的有限几何级数逆}{1-a的有限几何级数逆}}
\label{ex:1:38}

\paragraph{题目}
幺环中若$a^n=0$，则$a$称\textbf{幂零元}，且$1-a$可逆。

\paragraph{考点与思路}
构造有限几何级数，并分别核对左右乘积。

\paragraph{所用结论与原文依据}
\noteref{unit}；\noteref{ringlaws}。

% 原解来源：chapters/Part 2/chapter1-exercises.tex:387；第1章原\exref{1}{38}。

\begin{proof}
令$u=1+a+\cdots+a^{n-1}$。由分配律及幂法则，
\begin{align*}
(1-a)u&=(1+a+\cdots+a^{n-1})-(a+a^2+\cdots+a^n)\\
&=1-a^n=1,\\
u(1-a)&=(1+a+\cdots+a^{n-1})-(a+a^2+\cdots+a^n)\\
&=1-a^n=1.
\end{align*}
所以$u$是$1-a$的双侧逆元，即$(1-a)^{-1}=1+a+\cdots+a^{n-1}$。这里不要求$R$交换，因为计算中只出现同一个元素$a$的幂，它们彼此交换。
\end{proof}

\paragraph{复习提示}
非交换环中必须核对双侧逆元；同一元素的幂仍可按通常幂法则计算。


\subsection{习题39：\texorpdfstring{交换环的幂零元构成理想}{交换环的幂零元构成理想}}
\label{ex:1:39}

\paragraph{题目}
证明交换环中全体幂零元构成理想，并说明非交换环中该结论不成立。

\paragraph{考点与思路}
它是\exref{1}{34}中$I=\{0\}$的特例；复习时也应能独立写出幂次数的论证。

\paragraph{所用结论与原文依据}
\noteref{ideal}；\noteref{ringlaws}；\exref{1}{34}。补充前置工具：交换环的二项式展开。

% 原解来源：chapters/Part 2/chapter1-exercises.tex:395；第1章原\exref{1}{39}。

\begin{proof}
记$\mathcal N=\{a\in R:\text{存在正整数 }n\text{ 使 }a^n=0\}$。因为$0\in\mathcal N$，它非空。
若$a,b\in\mathcal N$，选$n,m\geq1$使$a^n=b^m=0$。若环无单位元，二项式展开的端项直接表示$(-b)^{n+m-1}$与$a^{n+m-1}$，不使用零次幂。展开$(a-b)^{n+m-1}$，其中任意一项$a^k b^{n+m-1-k}$都满足$k\geq n$或$n+m-1-k\geq m$，因而为零。所以
\[
(a-b)^{n+m-1}=0,
\]
即$a-b\in\mathcal N$。若$r\in R$，交换性给出$(ra)^n=r^na^n=0$，所以$ra=ar\in\mathcal N$。因此$\mathcal N$为理想。

非交换环的反例取$M_2(\Q)$中的
\[
A=\begin{pmatrix}0&1\\0&0\end{pmatrix},\qquad
B=\begin{pmatrix}0&0\\1&0\end{pmatrix}.
\]
有$A^2=B^2=0$，但$(A+B)^2=E$。若$A+B$幂零，某个正整数$k$应满足$(A+B)^k=0$，于是$(A+B)^{2k}=0$；另一方面$(A+B)^{2k}=E$，矛盾。因此幂零元之和未必幂零，幂零元集甚至不一定对加法封闭。
\end{proof}

\paragraph{复习提示}
证明中交换性的用途是二项式展开以及$(ra)^n=r^na^n$，不能把这两步用于任意非交换环。


\subsection{习题40：\texorpdfstring{有限幺环的单侧逆也是双侧逆}{有限幺环的单侧逆也是双侧逆}}
\label{ex:1:40}

\paragraph{题目}
若$R$有限且$xy=1$，则$yx=1$。

\paragraph{考点与思路}
将单侧逆转为有限集合上的单侧逆映射，再使用有限性。

\paragraph{所用结论与原文依据}
\noteref{unit}；\noteref{ringlaws}。

% 原解来源：chapters/Part 2/chapter1-exercises.tex:404；第1章原\exref{1}{40}。

\begin{proof}
在集合$R$上定义$L_x(r)=xr$、$L_y(r)=yr$。由结合律及$xy=1$，对每个$r\in R$都有
\[
L_xL_y(r)=x(yr)=(xy)r=r,
\]
故$L_xL_y=\mathrm{id}_R$，特别地$L_x$满射。由于$R$有限，$L_x$也单射。

于是对任意$r$，从
\[
L_x(L_yL_x(r))=(L_xL_y)L_x(r)=L_x(r)
\]
由$L_x$的单射性得到$L_yL_x(r)=r$。令$r=1$，便得$yx=1$。
\end{proof}

\paragraph{复习提示}
有限性用在满射推出单射，不能在无限环中照搬这一步。


\subsection{习题41：\texorpdfstring{单侧逆不成双侧逆时有无穷多个右逆}{单侧逆不成双侧逆时有无穷多个右逆}}
\label{ex:1:41}

\paragraph{题目}
若幺环中$ab=1$但$ba\neq1$，则满足$ax=1$的$x$有无穷多个。

\paragraph{考点与思路}
沿原解构造一列右逆，再用适当次幂右乘来区分它们。

\paragraph{所用结论与原文依据}
\noteref{unit}；\noteref{ringlaws}。

% 原解来源：chapters/Part 2/chapter1-exercises.tex:413；第1章原\exref{1}{41}。

\begin{proof}
令$c=1-ba$，则$c\ne0$，且
\[
ac=a-aba=0,\qquad cb=b-bab=0.
\]
对每个整数$n\geq0$定义$x_n=b+ca^n$。于是
\[
ax_n=ab+(ac)a^n=1,
\]
所以每个$x_n$都是$a$的右逆。

若存在$0\leq i<j$使$x_i=x_j$，则$ca^i=ca^j$。由$ab=1$，逐次消去相邻$ab$可得$a^jb^j=1$及$a^ib^j=b^{j-i}$。在上述等式右边乘$b^j$，得到
\[
cb^{j-i}=c.
\]
但$j-i\geq1$且$cb=0$，故左边为0，推出$c=0$，矛盾。因此$x_0,x_1,\ldots$两两不同，$a$有无穷多个右逆。
\end{proof}

\paragraph{复习提示}
不能只给出无穷多个公式，还须证明公式表示的元素两两不同。


\subsection{习题42：\texorpdfstring{唯一的广义逆推出体}{唯一的广义逆推出体}}
\label{ex:1:42}

\paragraph{题目}
环$R$至少有两个元素，若每个$a\neq0$恰有一个$b$满足$aba=a$，则：无零因子，$bab=b$，有单位元，且$R$为体。

\paragraph{考点与思路}
先由广义逆的唯一性排除零因子，再构造左右单位元。

\paragraph{所用结论与原文依据}
\noteref{zerodiv}；\noteref{unit}；\noteref{field}；\exref{1}{31}。

% 原解来源：chapters/Part 2/chapter1-exercises.tex:424；第1章原\exref{1}{42}。

\begin{proof}
\begin{enumerate}
\item 固定$a\ne0$，取唯一的$b$使$aba=a$。若$ax=0$，则
\[
a(b+x)a=aba+(ax)a=a.
\]
由$b$的唯一性，$b+x=b$，故$x=0$。若$xa=0$，同样有$a(b+x)a=a+a(xa)=a$，仍得$x=0$。这说明$R$无左右零因子。

\item 由$aba=a\ne0$知$b\ne0$，又
\[
a(bab-b)=(aba)b-ab=ab-ab=0.
\]
因为$a\ne0$且环无零因子，$bab=b$。

\item 任取$r\in R$，有
\[
a(bar-r)=(aba)r-ar=0,\qquad
(rab-r)a=r(aba)-ra=0.
\]
由无零因子得$bar=r$、$rab=r$。故$ba$是左单位元，$ab$是右单位元。第一章习题31第一项给出$ba=ab=1_R$，此元素是整个环的双侧单位元。

\item 对每个非零$a$重复上述构造，都得到双侧逆元$b$。单位元唯一，故各次构造得到的是同一个$1_R$。环至少有两个元素，故$1_R\ne0_R$；否则$r=1_Rr=0_R$对所有$r$成立，环只有一个元素。按体的定义，$R$为体。
\end{enumerate}
\end{proof}

\paragraph{复习提示}
单位元属于整个环，不能只证明它对所选的$a$起作用。


\subsection{习题43：\texorpdfstring{连续函数环}{连续函数环}}
\label{ex:1:43}

\paragraph{题目}
记$C[0,1]$为$[0,1]$上全部实值连续函数构成的环，加法和乘法逐点定义。
\begin{enumerate}
\item 每个真理想$I\ne C[0,1]$都存在公共零点，即存在$t_0\in[0,1]$使所有$f\in I$均满足$f(t_0)=0$。
\item 对非零函数$f\in C[0,1]$，证明$f$为零因子，当且仅当其零点集包含一个非空的相对开区间。
\end{enumerate}
相对开区间指$(\alpha,\beta)\cap[0,1]$形式的非空集合，其中$\alpha<\beta$为实数。这里“非零函数”指不恒为零的函数。

\paragraph{考点与思路}
第一问用紧致性选有限子覆盖，再以平方和构造处处非零的理想元素；第二问使用连续函数在非零点附近仍非零的性质。

\paragraph{所用结论与原文依据}
\noteref{ideal}；\noteref{zerodiv}；\noteref{idealunit}。补充前置工具：闭区间紧致性、连续函数的非零邻域及连续倒数，均为分析前置工具。

% 原解来源：chapters/Part 2/chapter1-exercises.tex:439；第1章原\exref{1}{43}。

\begin{proof}
\begin{enumerate}\item
反设$I$没有公共零点，则对每个$t\in[0,1]$，都可选$f_t\in I$使$f_t(t)\ne0$。
由连续性存在$t$的相对开邻域$U_t$，使$f_t$在$U_t$上处处非零。这些邻域覆盖$[0,1]$。

闭区间$[0,1]$紧致，故可选有限子覆盖$U_{t_1},\ldots,U_{t_m}$。记$f_j=f_{t_j}$，并令
\[
g=f_1^2+\cdots+f_m^2.
\]
因为$f_j\in I$，理想吸收性给出$f_j^2\in I$，加法封闭性给出$g\in I$。
任意$x\in[0,1]$至少属于某个$U_{t_j}$，因此对应$f_j(x)\ne0$，从而$g(x)=\sum_{j=1}^m f_j(x)^2>0$。
所以函数$h(x)=1/g(x)$在整个区间连续，属于$C[0,1]$。

由理想吸收性，常值函数$1=gh\in I$。于是任取$u\in C[0,1]$都有$u=u\cdot1\in I$，这说明$I=C[0,1]$，与$I$为真理想矛盾。故必有公共零点。

\item\textbf{必要性。}
设$f\ne0$为零因子，则存在不恒为零的$g\in C[0,1]$使$fg=0$。选$t\in[0,1]$满足$g(t)\ne0$。
连续性给出某个$\delta>0$，使得当$x\in[0,1]$且$|x-t|<\delta$时，
\[
|g(x)-g(t)|<\frac{|g(t)|}{2},\qquad
|g(x)|>\frac{|g(t)|}{2}>0.
\]
在非空相对开区间$(t-\delta,t+\delta)\cap[0,1]$上，$f(x)g(x)=0$且$g(x)\ne0$，所以$f(x)=0$。因此$f$的零点集包含这样的区间。

\textbf{充分性。}
反过来，设$f$在非空相对开区间$J$上恒为零。这样的$J$含有一个位于$(0,1)$内的小闭区间$[c-\varepsilon,c+\varepsilon]$，其中$\varepsilon>0$。这是因为非空相对开集即使包含端点，也包含相邻的内部点，故可在其内部选$c$及足够小的$\varepsilon$。
定义连续的三角形函数
\[
g(x)=\max\left\{0,1-\frac{|x-c|}{\varepsilon}\right\},\qquad x\in[0,1].
\]
它满足$g(c)=1$，所以不恒为零；在$[c-\varepsilon,c+\varepsilon]$外为零，而在该闭区间内$f$为零。因此对所有$x\in[0,1]$均有$f(x)g(x)=0$，即$fg=0$。由于$f\ne0$，$f$为零因子。
\end{enumerate}
\end{proof}

\paragraph{复习提示}
“每个$f\in I$各有零点”与“所有$f\in I$有同一个公共零点”的量词次序不同。第二问不能只要求$f$有一个零点；必须有一段区间。零函数满足零点集条件，但按约定不称零因子，所以题目须限定$f\ne0$。


\subsection{习题44：\texorpdfstring{有限域上的矩阵计数}{有限域上的矩阵计数}}
\label{ex:1:44}

\paragraph{题目}
设$p$为素数，$F=\F_p$为$p$元域，$n\geq1$。证明
\[
|M_n(F)|=p^{n^2},\qquad
|\GL_n(F)|=\prod_{k=0}^{n-1}(p^n-p^k).
\]
这里$M_n(F)$是全体$n$阶矩阵的集合，$\GL_n(F)$是其中的可逆矩阵乘法群。

\paragraph{考点与思路}
可逆矩阵的列向量恰为$F^n$的一组有序基。逐列选取时，排除前面列向量张成的子空间。

\paragraph{所用结论与原文依据}
\noteref{iso}。补充前置工具：线性代数中可逆矩阵与有序基的等价判别，以及有限集的乘法计数原理。

% 原解来源：chapters/Part 2/chapter1-exercises.tex:453；第1章原\exref{1}{44}。

\begin{proof}
一个$n$阶矩阵共有$n^2$个位置，每一位置均可独立选取$F$的$p$个元素之一。因此乘法计数原理给出$|M_n(F)|=p^{n^2}$。

现在计数可逆矩阵。线性代数的可逆性判别说明，矩阵可逆当且仅当其$n$个列向量线性无关。第一列不能是零向量，所以有$p^n-1$种选择。

假设已选出线性无关的前$k$列$v_1,\ldots,v_k$，其中$0\leq k<n$。其张成空间中的向量均可写成
\[
\lambda_1v_1+\cdots+\lambda_kv_k,\qquad \lambda_i\in F.
\]
因为这$k$列线性无关，不同系数组给出不同向量：若两组给出相同向量，相减后得到线性关系，所有系数差必须为零。因此该张成空间恰有$p^k$个向量。
第$k+1$列必须位于该空间之外，才能保持列向量线性无关，所以有$p^n-p^k$种选择。

依次选满$n$列，每个有序选择恰对应一个矩阵，没有额外的排列除数。故
\[
|\GL_n(F)|=(p^n-1)(p^n-p)\cdots(p^n-p^{n-1}).
\]
这正是所求乘积。
\end{proof}

\paragraph{复习提示}
当$n=1$时结果为$p-1$，对应域中非零元素的个数。关键区别是：每一列都非零不足以保证可逆，还必须排除前面列向量的整个张成空间。


\subsection{习题45：\texorpdfstring{华罗庚恒等式}{华罗庚恒等式}}
\label{ex:1:45}

\paragraph{题目}
在体$K$中，若$a,b\neq0$且$ab\neq1$，则
\[
a-\bigl(a^{-1}+(b^{-1}-a)^{-1}\bigr)^{-1}=aba.
\]

\paragraph{考点与思路}
先确认嵌套逆元全部存在，再沿原解选择$d=a-aba$并核对左乘积。

\paragraph{所用结论与原文依据}
\noteref{field}；\noteref{unit}；\noteref{ringlaws}。

% 原解来源：chapters/Part 2/chapter1-exercises.tex:467；第1章原\exref{1}{45}。

\begin{proof}
令$c=b^{-1}-a$。若$c=0$，则$a=b^{-1}$，所以$ab=1$，与题设矛盾。故$c\ne0$，$c^{-1}$存在。

令$d=cba$。由$a,b,c$均非零且体无零因子，$d\ne0$，所以$d$可逆。展开得
\[
d=(b^{-1}-a)ba=a-aba.
\]
保留乘法次序，计算
\begin{align*}
(a^{-1}+c^{-1})d
&=a^{-1}(a-aba)+c^{-1}(cba)\\
&=(1-ba)+ba=1.
\end{align*}
右乘$d^{-1}$，得$a^{-1}+c^{-1}=d^{-1}$。所以$a^{-1}+c^{-1}$确实非零且可逆，其逆为$d$。从而
\[
a-\bigl(a^{-1}+(b^{-1}-a)^{-1}\bigr)^{-1}
=a-d=aba.
\]
原笔记还给出辅助恒等式。对$x\ne0,1$，由$x^{-1}-1=x^{-1}(1-x)$及乘积取逆反序公式，
\begin{align*}
(x^{-1}-1)^{-1}&=(1-x)^{-1}x\\
&=(1-x)^{-1}\bigl(1-(1-x)\bigr)\\
&=(1-x)^{-1}-1.
\end{align*}
这些计算均未交换不同元素的乘法次序。
\end{proof}

\paragraph{复习提示}
嵌套表达式每一层取逆都必须先说明对象非零，不能只在形式上变形。


